% ---------------------------------------------------------------------
% ---------------------------------------------------------------------

\IEEEPARstart{T}{he} quality of AI-generated speech has crossed a perceptual
threshold. Neural codec language models such as VALL-E~\cite{wang2023valle},
diffusion- and flow-matching-based synthesizers such as StyleTTS~2~\cite{li2023styletts2}
and F5-TTS~\cite{chen2024f5tts}, and massively multilingual zero-shot systems such as
XTTS~\cite{casanova2024xtts}, Seed-TTS~\cite{anastassiou2024seedtts}, and
CosyVoice~\cite{du2024cosyvoice} can now clone a speaker's voice from a few seconds of
reference audio and synthesize speech that human listeners and conventional classifiers
cannot reliably distinguish from genuine recordings. While these systems enable valuable
applications in accessibility, dubbing, and creative production, they have also driven a
sharp rise in documented harms: voice-cloning fraud against individuals and financial
institutions, fabricated audio of political figures circulated during elections, and
evasion of voice-biometric authentication. An analysis of deepfakes actually circulated
online during 2024 found that audio deepfakes are among the fastest-growing and least
reliably detected modalities~\cite{chandra2025deepfakeeval}, and large-scale fake-speech
corpora are now being generated with large language models in the loop to simulate
disinformation campaigns~\cite{luong2024llamapartialspoof}.

The anti-spoofing community has responded with increasingly capable countermeasures.
On the standard ASVspoof~2019 Logical Access (LA) benchmark~\cite{wang2020asvspoof2019},
systems built on self-supervised learning (SSL) speech foundation models with
graph-attention back-ends~\cite{tak2022automatic,jung2022aasist} achieve equal error
rates (EER) below 1\%, and recent state-space-model
detectors~\cite{xiao2025xlsrmamba,chen2024rawbmamba} and lightweight foundation-model
back-ends~\cite{liu2025nes2net} have pushed in-domain accuracy further still. Yet a
growing body of evidence---from the ASVspoof~5 challenge~\cite{wang2024asvspoof5,
wang2025asvspoof5journal} to systematic cross-domain
studies~\cite{muller2022inthewild,muller2024harder}---shows that these in-domain numbers
are a poor predictor of real-world reliability. We identify three failure modes that
jointly block deployment.

\textbf{G1---Generalization to unseen generators.} A detector trained on a fixed set of
text-to-speech (TTS) and voice-conversion (VC) algorithms sees its EER degrade from
below 1\% to 5--30\% when evaluated on synthesizers absent from its training
corpus~\cite{muller2022inthewild,muller2024harder}. Because new synthesis families
appear monthly---neural codec language models and flow-matching synthesizers did not
exist when the dominant training corpora were collected---a practical detector must
generalize \emph{zero-shot}, without retraining, to generator families it has never
seen. The Codecfake study~\cite{xie2025codecfake} demonstrates that detectors trained on
conventional vocoder artifacts largely fail on speech produced through neural audio
codecs, precisely the mechanism used by the newest generation of synthesizers.

\textbf{G2---Channel robustness.} Real-world audio reaches a detector after lossy
compression (MP3, AAC, Opus), telephony coding (AMR, G.711), reverberation, and additive
noise---conditions imposed by messaging platforms, social media transcoding, and phone
networks. The ASVspoof~2021 evaluation~\cite{liu2023asvspoof2021} showed that codec and
compression variability alone can multiply EER several-fold, and the interaction of
channel distortion with spoofing artifacts remains poorly understood. Most published
evaluations still use clean, studio-quality audio.

\textbf{G3---Calibration and explainability.} Detection scores feed downstream
decisions---content moderation, forensic analysis, legal evidence---that require
\emph{calibrated} confidence estimates, not merely a ranking. Recent work shows that
modern SSL-based detectors are systematically over-confident, particularly under domain
shift~\cite{pascu2024calibrated}. Forensic practice additionally demands interpretable
evidence for individual decisions, which binary scores do not provide.

Recent large audio-language models (ALLMs) such as Qwen2-Audio~\cite{chu2024qwen2audio}
and SALMONN~\cite{tang2024salmonn} raise a natural question: does the general-purpose
audio understanding of foundation models subsume the deepfake-detection task? Early
evidence suggests not---prompted zero-shot, ALLMs perform near chance at distinguishing
bona-fide from synthetic speech~\cite{gong2025allm4add}. This stems from their objective:
general audio models compress acoustic semantics while abstracting away the low-level
phase jitter and vocoder artifacts essential for forensic verification. A comprehensive
study placing speech foundation representations alongside multiple generations of
specialized countermeasures under controlled cross-domain, channel-robustness, and
calibration protocols remains an urgent priority. We undertake this evaluation here.

This paper reframes synthetic-speech detection from closed-set binary classification to
\emph{open-set, channel-robust, calibrated} detection, and proposes
\textbf{AuralGuard}, a multi-view one-class detection framework addressing G1--G3
jointly. AuralGuard rests on four design decisions. \emph{First}, it fuses two
complementary views of the input: View~A, a self-supervised foundation-model
representation (WavLM-Large~\cite{chen2022wavlm}) with \emph{learnable layer attention}
that adaptively selects the transformer depths most informative for artifact detection;
and View~B, a hand-crafted feature stream that explicitly exposes low-level vocoder
artifacts through linear-frequency cepstral coefficients (LFCC), modified group delay,
and constant-Q phase---cues that SSL models, pre-trained for semantic tasks, are prone
to discard. \emph{Second}, a gated cross-attention module fuses the two views before a
spectro-temporal graph-attention back-end~\cite{jung2022aasist} produces an
utterance-level embedding. \emph{Third}, a one-class contrastive objective combining
OC-Softmax~\cite{zhang2021oneclass} with a supervised-contrastive
auxiliary~\cite{khosla2020supcon} learns a compact decision region around bona-fide
speech, so that unseen spoofing algorithms are rejected as out-of-distribution rather
than mapped onto decision boundaries fit to known attacks. \emph{Fourth}, a
channel-robustness augmentation curriculum---RawBoost~\cite{tak2022rawboost}, a
realistic codec chain (MP3, AAC, Opus, AMR, G.711), measured room impulse responses,
and additive noise---aligns the training distribution with deployment channels.

We evaluate AuralGuard under a strict train-once, evaluate-everywhere protocol: all models are trained exclusively on the ASVspoof~2019~LA training partition and evaluated without fine-tuning or adaptation on held-out in-domain (ASVspoof~2019~LA eval) and challenging out-of-domain zero-shot corpora: In-the-Wild~\cite{muller2022inthewild} (probing real-world web audio and transmission noise) and WaveFake~\cite{frank2021wavefake} (probing unseen neural vocoders). Baselines span four generations of countermeasure design: hand-crafted features with light CNNs (LFCC-LCNN)~\cite{lavrentyeva2019stc}, end-to-end raw-waveform networks (RawNet2)~\cite{tak2021rawnet2}, graph-attention architectures (AASIST)~\cite{jung2022aasist}, and SSL-based foundation-model detectors with one-class scoring (WavLM+AASIST+OCS)~\cite{chen2022wavlm,zhang2021oneclass}. Beyond primary EER and minimum tandem detection cost (min~t-DCF)~\cite{kinnunen2020tdcf}, we report 95\% bootstrap confidence intervals, AUROC, balanced accuracy, F1-scores, expected calibration error (ECE), and Brier scores to assess forensic reliability under distribution shift~\cite{oyucu2026generalization,nagar2026hybrid}.

\noindent\textbf{Research Questions.} To systematically investigate the limits of generalizability and forensic calibration in synthetic speech detection, this study is structured around four primary research questions:
\begin{enumerate}
    \item \textbf{RQ1 (Cross-Corpus Generalization):} Can a multi-view countermeasure trained strictly on clean in-domain speech (ASVspoof~2019~LA) generalize zero-shot to unconstrained web audio (In-the-Wild) and multi-vocoder speech (WaveFake) without catastrophic performance collapse?
    \item \textbf{RQ2 (Forensic Probability Calibration):} Does an open-set one-class hyperspherical objective (OC-Softmax with SupCon regularization) produce well-calibrated posterior probability estimates and bounded confidence intervals under severe acoustic distribution shift, compared to standard binary cross-entropy and legacy acoustic countermeasures?
    \item \textbf{RQ3 (Multi-View Feature Complementarity):} What are the individual and synergistic contributions of the self-supervised foundation-model front-end (View~A), the physical phase-artifact stream (View~B), and bidirectional gated cross-attention fusion when exposed to lossy transmission codecs and diverse vocoder algorithms?
    \item \textbf{RQ4 (Parameter Efficiency and Operational Feasibility):} Can competitive zero-shot discrimination and forensic calibration be achieved while keeping the 315M foundation backbone frozen, and does the resulting inference latency support real-time operation on commodity hardware?
\end{enumerate}

The contributions of this work are as follows:
\begin{enumerate}
    \item \textbf{Multi-View Architecture.} AuralGuard, a novel detection framework that unifies an adaptive SSL foundation-model stream (WavLM-Large with learnable layer attention and entropy regularization) with an explicit low-level vocoder-artifact branch (LFCC, modified group delay, and CQT phase) via gated cross-attention fusion and a spectro-temporal graph back-end (Section~\ref{sec:method}).
    \item \textbf{One-Class Optimization.} An Orthogonal Centroid Scoring (OCS) one-class training objective that tightly bounds the bona-fide speech manifold, enabling robust rejection of unseen synthesis artifacts as out-of-distribution outliers without overfitting to training-set artifacts (Section~\ref{ssec:objective}).
    \item \textbf{Comprehensive Generational Benchmark.} A controlled empirical comparison of six models spanning four architectural generations under identical training data, feature extractors, and evaluation protocols across in-domain and zero-shot cross-corpus shifts (Section~\ref{ssec:generalization}).
    \item \textbf{Superior Probability Calibration.} Demonstration that multi-view one-class learning yields superior probability calibration under severe domain shift, achieving an In-the-Wild ECE of 0.0965 (v1) and 0.1716 (v2) compared to 0.2320 for single-view WavLM+OCS and $>$0.55 for legacy end-to-end networks, providing trustworthy reliability for forensic applications (Section~\ref{ssec:calibration}).
    \item \textbf{Open Reproducibility.} Complete source code, configuration files, trained model checkpoints, and evaluation artifacts are publicly released at GitHub (\url{https://github.com/MIHMahmudEli/auralguardv2}) and Hugging Face (\url{https://huggingface.co/MoshinAli/auralguard-checkpoints}) to support reproducible research in synthetic speech forensics.
\end{enumerate}

The remainder of the paper is organized as follows. Section~\ref{sec:related} reviews countermeasure architectures, speech foundation models, one-class learning, and recent cross-domain generalization studies. Section~\ref{sec:method} details the AuralGuard architecture, dual-branch feature representations, fusion, and optimization objective. Section~\ref{sec:setup} describes datasets, baselines, implementation details, and evaluation metrics. Section~\ref{sec:results} presents in-domain performance, zero-shot generalization, calibration analysis, and computational efficiency. Section~\ref{sec:discussion} discusses practical implications, limitations, and ethical considerations, and Section~\ref{sec:conclusion} concludes.
